% ============================================================================
\chapter{Ступень IV. Элементарные частицы, ядра, гравитация и космология}
\label{ch:course-stage4-particles-cosmos}
% ============================================================================

\section{Урок 12. Запуск Вселенной из слова \texorpdfstring{$\mathrm{IV} = 263 = 407_8$}{IV = 263 = 407\_8}: тишина октанта, храповик \texorpdfstring{$-1$}{-1} и рождение вихрей при \texorpdfstring{$t = 2$}{t = 2}}

Запустим теперь наш полный 3D клеточный автомат из канонического начального состояния $\Psi^{(0)} = \mathrm{IV}(3) = 263_{10} = 407_8$ в шаре радиуса $|\square(3)| = 6$:
\begin{enumerate}[label=\textbf{(\alph*)},noitemsep]
\item \textbf{Тишина внутренних октантов:} внутри однородного октанта в узле $(2,2,2)$ сырой фазовый импульс равен $\Phi_{\mathrm{raw}} = +15$, но поскольку $15 < 41 = \Delta\varphi_{\mathrm{disc}}$, мёртвая зона Гейзенберга гасит его в \textbf{точный ноль $\Phi_{\mathrm{gate}}(2,2,2) = 0$} (внутренность октанта является неподвижной точкой гомеостаза!).
\item \textbf{Киральный квантовый храповик $-1$ тик на доменной стенке:} на границе октантов $z = 0 / -1$ в паре узлов $(0,2,2)$ и $(-1,2,2)$ чётные голономии совпадают, нечётные токи взаимно компенсируются, но за счёт точного целочисленного деления с округлением вниз ($\lfloor -N/D\rfloor = -\lfloor N/D\rfloor - 1$ при ненулевом остатке) импульсы ворот равны $+76$ и $-77$ тиков, давая \textbf{киральную асимметрию $(+76) + (-77) = -1$ тик}! (Если заменить деление с округлением вниз на симметричное усечение к нулю, эта асимметрия между материей и антиматерией исчезает).
\item \textbf{Рождение вихревой материи на такте $t = 2$:} на центральной плакете $(-1,-1,0)$ вихревое число Черна равно $n_f(t=0) = 0$, $n_f(t=1) = 0$ и на втором такте скачком переходит в $\mathbf{n_f(t=2) = +1}$ с квантом потока Максвелла $B_f = \mathbf{512}$, порождая во всём 3D-шаре ровно $\mathbf{86 = (38, 16, 32)}$ вихревых плакет (и барионную асимметрию $\eta_B$ без подгоночных параметров!).
\end{enumerate}

\section{Урок 13. Симплициальный атлас \texorpdfstring{$\mathcal{C}_{13}$}{C\_13}, расшифровка «магических букв» групп симметрий, постоянная тонкой структуры \texorpdfstring{$\alpha_{\mathrm{fs}}^{-1} = 137{,}035999\dots$}{alpha\^-1 = 137.035999...} и спектр частиц}

Прежде чем мы выведем калибровочную группу Стандартной модели $\mathrm{U}(1)\times\mathrm{SU}(2)\times\mathrm{SU}(3)$, снимем завесу тайны с того, почему математики и физики обозначают группы симметрий такими «магическими» наборами латинских и готических букв.

\begin{quote}
\small\textbf{Историческая и терминологическая справка: Расшифровка «магических букв» групп симметрий (Эварист Галуа, 1832; Софус Ли, 1870-е; Герман Вейль, 1929).}
В ночь перед роковой дуэлью в мае 1832~года 20-летний французский математик Эварист Галуа ввёл понятие \emph{группы} перестановок симметрии. В 1870-х годах норвежский математик Софус Ли построил теорию непрерывных групп вращений, а в XX~веке за ними закрепились короткие английские аббревиатуры, которые в учебниках почти никогда не переводят на человеческий язык. Вот их полный словарь:
\begin{itemize}[noitemsep,leftmargin=*]
\item \textbf{В скобках $(n)$ или $(\mathbb{Z}), (\mathbb{R}), (\mathbb{C})$} всегда указывается размер таблицы (матрицы $n\times n$) и какие числа стоят внутри: целые~$\mathbb{Z}$, вещественные~$\mathbb{R}$ или комплексные~$\mathbb{C}$.
\item \textbf{Буква \texttt{L} = Linear («Линейная»)} и \textbf{\texttt{GL}$_n$ = General Linear («Общая линейная группа»):} вообще все обратимые матрицы $n\times n$ с ненулевым определителем ($\det \ne 0$). Они могут как угодно крутить, растягивать и отражать пространство, лишь бы не сплющивать объём в ноль.
\item \textbf{Буква \texttt{S} в самом начале = Special («Специальная», то есть $\det = +1$):} как только перед любой группой ставят букву~\texttt{S}, это означает одно простое правило --- \textbf{определитель матрицы строго равен $+1$} (преобразование сохраняет элементарный объём ячейки и не выворачивает правую руку в левую). Поэтому \textbf{$\mathrm{SL}_2(\mathbb{Z})$ = Special Linear $2\times 2$ over $\mathbb{Z}$} --- это целочисленные матрицы $\begin{pmatrix} a & b \\ c & d \end{pmatrix}$ с единичной площадью клетки $ad - bc = 1$ (наш шаг Верле и каскад Паде--Минковского!).
\item \textbf{Буква \texttt{O} = Orthogonal («Ортогональная» --- повороты в обычном вещественном пространстве):} матрицы, сохраняющие евклидову длину вектора $x^2 + y^2 + z^2 = R^2$ ($\det = \pm 1$). Группа \textbf{$\mathrm{O}(3)$} включает все повороты и зеркальные отражения 3D-шара (а её конечная подгруппа \textbf{$O_h$} из $48$ элементов --- все повороты и отражения нашего куба и кубооктаэдра). Если же приписать спереди букву~\texttt{S}, получим \textbf{$\mathrm{SO}(3)$ = Special Orthogonal} --- чистые трёхмерные повороты твёрдого тела без зеркал ($\det = +1$), которые задаются тремя углами Эйлера в трёх плоскостях $(xy, xz, yz)$ кубооктаэдра.
\item \textbf{Буква \texttt{U} = Unitary («Унитарная» --- повороты комплексных волн с фазой $e^{i\varphi}$):} когда у ячейки есть фаза на окружности, роль длины вектора играет сумма квадратов модулей $|\psi_1|^2 + \dots + |\psi_n|^2$. Преобразования, сохраняющие эту сумму, называются \emph{унитарными}:
  \begin{itemize}[noitemsep]
  \item \textbf{$\mathrm{U}(1)$} ($1\times 1$, размерность $1$) --- обычный сдвиг одной фазы $\varphi \mapsto \varphi + \Delta\varphi$ по единичной окружности (наше кольцо $\mathbb{Z}_{512}$; в физике это \textbf{электромагнетизм}, введённый Германом Вейлем в 1929~году);
  \item \textbf{$\mathrm{SU}(2)$} ($2\times 2$ с $\det = +1$, размерность $2^2 - 1 = \mathbf{3}$) --- комплексные повороты двухкомпонентного спинора $(\psi_0, \psi_1)$ по $d = 3$ пространственным осям (отвечает за спин электрона и \textbf{слабое взаимодействие});
  \item \textbf{$\mathrm{SU}(3)$} ($3\times 3$ с $\det = +1$, размерность $3^2 - 1 = \mathbf{8}$) --- комплексные перестановки и повороты между тремя плоскостями/цветами на $8$ треугольных гранях кубооктаэдра (отвечает за \textbf{сильное ядерное взаимодействие} кварков и глюонов).
  \end{itemize}
\item \textbf{Буква \texttt{P} в самом начале = Projective («Проективная», например $\mathrm{PGL}_2, \mathrm{PSL}_2$):} означает, что матрица $\begin{pmatrix} a & b \\ c & d \end{pmatrix}$ действует на дробь $w \mapsto \frac{aw+b}{cw+d}$ на проективной сфере Римана $\mathbb{P}^1$, где общий множитель числителя и знаменателя сокращается.
\item \textbf{Маленькие готические буквы $\mathfrak{u}(1), \mathfrak{su}(2), \mathfrak{su}(3)$ (алгебры Ли):} большие латинские буквы ($\mathrm{SU}(3)$) обозначают уже готовый конечный поворот, а маленькие готические буквы ($\mathfrak{su}(3)$) --- мгновенную угловую скорость (бесконечно малый толчок или коммутатор $[T_a, T_b] = T_a T_b - T_b T_a$) в начале этого поворота.
\end{itemize}
\end{quote}

Теперь посчитаем все полные подграфы (симплициальные клики $K_m$, где каждые две вершины соединены ребром длины $\|u-v\|_2^2 = 2$) внутри 13-узлового одетого кубооктаэдра $\mathcal{C}_{13} = \{(0,0,0)\}\cup O_2$:
\begin{itemize}[noitemsep]
\item \textbf{$K_1$ (вершины):} $1\text{ (центр)} + 12\text{ (оболочка }O_2\text{)} = \mathbf{13 = N_\varphi(3)}$;
\item \textbf{$K_2$ (рёбра):} $12\text{ радиальных} + 24\text{ поверхностных} = \mathbf{36 = |\square(3)|^2 = B_1(3)}$;
\item \textbf{$K_3$ (треугольники):} $24\text{ радиально-рёберных} + 8\text{ поверхностных} = \mathbf{32 = K_P(3)/2}$;
\item \textbf{$K_4$ (тетраэдры) и почему $K_{\ge 5} = 0$:} каждый из 8 поверхностных треугольников образует с центром $(0,0,0)$ правильный тетраэдр~$K_4$. А может ли на самой поверхности кубооктаэдра $O_2$ лежать четвёрка взаимно смежных вершин? Нет: окрестность любой вершины на поверхности $O_2$ состоит из $4$ вершин, образующих цикл $C_4$ из двух противоположных треугольников и двух квадратов, в котором нет ни одного треугольника $K_3$! Поэтому все 4-клики обязаны содержать центр $(0,0,0)$ и один из $8$ поверхностных треугольников ($\mathbf{K_4 = 8 = |{\triangle}(3)|}$), а кликов из $5$ и более вершин на $\mathcal{C}_{13}$ \textbf{не существует вовсе ($K_{\ge 5} \equiv 0$)}!
\end{itemize}
Отсюда:
\begin{enumerate}[label=\textbf{\arabic*.},noitemsep]
\item Симплициальная характеристика Эйлера равна $\chi(\mathcal{C}_{13}) = 13 - 36 + 32 - 8 = \mathbf{1 = \mu_{\mathrm{vac}}(3)}$.
\item Число независимых генераторов калибровочной группы Стандартной модели $\mathrm{U}(1)_Y\times\mathrm{SU}(2)_L\times\mathrm{SU}(3)_c$ (алгебры Ли $\mathfrak{u}(1)\oplus\mathfrak{su}(2)\oplus\mathfrak{su}(3)$) равно:
  \[
    \mu_{\mathrm{vac}}(3) + d + K_4(\mathcal{C}_{13}) \;=\; 1 + 3 + 8 \;=\; \mathbf{12 \;=\; \geomNeighborTwelve(3)}.
  \]
\item Относительно любой выбранной вершины $v_0\in O_2$ (ротор электрона) 13 узлов $\mathcal{C}_{13}$ делятся по квадрату расстояния до $v_0$ на четыре слоя $\mathbf{1\text{ (центр)} + 1\text{ ($v_0$)} + 4\text{ (смежные на $O_2$)} + 7\text{ (несмежные на $O_2$)} = 13}$, откуда сразу следуют углы смешивания нейтрино PMNS:
  \[
    \sin^2\theta_{12} \;=\; \frac{4}{13}, \qquad \sin^2\theta_{23} \;=\; \frac{7}{13}, \qquad \sin^2\theta_{13} \;=\; \frac{2}{92} \;=\; \frac{1}{46},
  \]
  а также древесный угол Вайнберга $\sin^2\theta_W^{(0)} = d / N_\varphi(3) = 3/13$ и однопетлевой одетый угол Вайнберга на масштабе $Z$-бозона $\sin^2\theta_W(M_Z) = \frac{3}{13} + \frac{\alpha_{\mathrm{fs}}}{16} \approx \mathbf{0{,}231225}$ (PDG: $0{,}23122(4)$).
\item \textbf{Гауссово кулоновское простое $\hat{M}(3) = 97$ и постоянная тонкой структуры $\alpha_{\mathrm{fs}}^{-1}$:}
  единственное размерностное совпадение $1 + \geomNeighborTwelve(d)\cdot |{\triangle}(d)| = d_{\mathrm{in}}(\square)^2 + \mathrm{mod\_bits}(d)^2$ при $d=3$ даёт норму гауссова простого $4 + 9i\in\mathbb{Z}[i]$:
  \[
    \hat{M}(3) \;=\; 1 + 12\times 8 \;=\; 4^2 + 9^2 \;=\; \mathbf{97}
    \qquad \bigl(\text{причём } (m_u/m_e)^2 + (m_d/m_e)^2 = 4^2 + 9^2 = 97!\bigr),
  \]
  а добавление 40 ненулевых мод порога Гейзенберга ($\Delta\varphi_{\mathrm{disc}} - 1 = 5\times 8 = 40$) даёт целую базу постоянной тонкой структуры $\lfloor\alpha_{\mathrm{fs}}^{-1}\rfloor = 97 + 40 = \mathbf{137}$. Поскольку симплициальный комплекс $\mathcal{C}_{13}$ обрывается строго на $K_4 = 8$ ($K_{\ge 5} \equiv 0$), учёт вкладов всех четырёх непустых симплексов $K_1, K_2, K_3, K_4$ на длине ребра кубооктаэдра $\|e\|_2 = \sqrt{d-1} = \sqrt{2}$ приводит к замкнутой алгебраической формуле в поле $\mathbb{Q}(\sqrt{2})$ без участия трансцендентного числа~$\pi$:
  \begin{equation}
    \label{eq:course-alpha-fs}
    \alpha_{\mathrm{fs}}^{-1}(3)
    \;=\; \frac{A_{\mathrm{den}}(3) + B_{\mathrm{den}}(3)\sqrt{2}}{A_{\mathrm{num}}(3) + B_{\mathrm{num}}(3)\sqrt{2}}
    \;=\; \frac{17\,351\,739\,044 - 679\sqrt{2}}{131\,726 + 89\,441\,954\sqrt{2}}
    \;=\; \mathbf{137{,}035999187\dots}
  \end{equation}
  которая попадает в эксперимент CODATA-2022 ($137{,}035999177(21)$) с точностью $\mathbf{0{,}47\sigma}$ без единого подгоночного параметра!
\item Магнитные моменты протона и нейтрона равны $\mu_p/\mu_N = 3 - 5/24 = \mathbf{67/24 = +2{,}7917}$ и $\mu_n/\mu_N = -2 + 2/24 = \mathbf{-23/12 = -1{,}9167}$ с изоскалярной суммой $67/24 - 46/24 = 21/24 = \mathbf{7/8}$, аксиальная константа равна $g_A = 5/4 + 1/39 = \mathbf{199/156 = 1{,}27564}$, а разность масс нейтрона и протона равна $(m_n - m_p)/m_e = 2|\mathcal{E}_{\mathrm{Heis}}| / K_P = 162 / 64 = \mathbf{81/32 = 2{,}53125}$ (PDG: $2{,}53099$).
\end{enumerate}

\section{Урок 14. Гравитация без сингулярности (\texorpdfstring{$833:1$}{833:1}), ядерная паста и гало вакуума \texorpdfstring{$16/3$}{16/3} без тёмной материи}

Откуда в дискретной решётке берётся гравитация? Посмотри на знаменатель импеданса ворот ячейки:
\[
  D(x) \;=\; \max\Bigl(\varepsilon,\; \varepsilon + 128\,\mu(x)^2 + \mu(x)\sum_{j=1}^{12}\mu(x+e_j)\,\mathcal{C}_3(\Delta\varphi_j)\Bigr),
  \qquad \varepsilon \;=\; K_P\cdot Q_\square \;=\; 64\times 128 \;=\; \mathbf{8\,192}.
\]
Когда в области пространства сосредоточена массивная материя ($\mu > 1$), знаменатель импеданса $D(x)$ возрастает, \textbf{замедляя фазовую скорость такта} $\Delta\varphi \propto 1/D(x)$ точно так же, как оптический показатель преломления $n(x)$ поворачивает световой луч в линзе!
\begin{itemize}[leftmargin=*]
\item \textbf{Отсутствие сингулярности чёрной дыры и предел сжатия $833:1$:}
  в уравнениях ОТО Эйнштейна при $r\to 0$ возникает деление на ноль ($1/0 = \infty$). В решётке кубооктаэдра даже при максимальной плотности $\mu = K_P = 64$ и противофазном горизонте Паули $\Delta\varphi = \Phi_{\mathrm{Pauli}} = 256$ ($\mathcal{C}_3(256) = -128$) знаменатель импеданса никогда не обращается в нуль, а ложится на квантовый пол $D_{\min} = \varepsilon = \mathbf{8\,192 > 0}$! При этом радиальный ток переноса на горизонте Паули тождественно замирает ($\mathcal{S}_3(256) = 0$, останавливая коллапс в фазовом кристалле Паули), отношение ненапряжённого импеданса стазиса ($\Delta\varphi=0$, $D_{\max} = 8\,192\times(1 + 13\times 64) = \mathbf{6\,823\,936}$) к импедансу горизонта равно точному целому числу:
  \[
    \frac{D_{\max}}{D_{\min}} \;=\; 1 + N_\varphi(3)\,K_P(3) \;=\; 1 + 13\times 64 \;=\; \mathbf{833 \;=\; 7^2 \times 17},
  \]
  а побитовая обратимость шага Верле сохраняет всю информацию без потерь (закрывая информационный парадокс Хокинга).
\item \textbf{Гало голономии вакуума $16/3$ без тёмной материи:}
  внешний вакуумный пол $\mu_{\mathrm{vac}} = 1$ вокруг галактик не пуст: отношение ёмкости $2|{\triangle}(3)| = 16$ спинорных каналов квадратной грани гало к $d = 3$ пространственным каналам барионного ядра равно:
  \[
    \frac{\Omega_{\mathrm{halo}}}{\Omega_b} \;=\; \frac{2\,|{\triangle}(3)|}{d} \;=\; \mathbf{\frac{16}{3} \;\approx\; 5{,}333},
    \qquad
    \frac{\Omega_{\mathrm{total}}}{\Omega_b} \;=\; \frac{16 + 3}{3} \;=\; \mathbf{\frac{19}{3}} \quad \bigl(19 = 3^3 - 2^3 = N_\varphi + |\square|\bigr),
  \]
  что воспроизводит наблюдаемое отношение Planck~2018 ($5{,}36\pm 0{,}07$) с тождественно нулевым остатком «тёмной материи» ($D_{\mathrm{dark}}\equiv 0$).
\end{itemize}

\section{Урок 15. Секунда без атома цезия, полупериодное горлышко \texorpdfstring{$\tau_{\mathrm{neck}} = 256$}{t = 256} и почему Вселенной не \texorpdfstring{$13{,}8$}{13.8}, а \texorpdfstring{$\sim 27{,}5\text{ млрд лет}$}{27.5 млрд лет} (разгадка ранних галактик JWST)}

В системе СИ секунда определена через $9\,192\,631\,770$ периодов излучения атома цезия-$133$. Но у самого носителя есть собственный фундаментальный такт! За $16 = d_{\mathrm{in}}(\square)^2 = 2|{\triangle}(3)|$ каскадных оборотов 9-битного кольца $\mathbb{Z}_{512}$ набирается $512^{16} = 2^{144} = 2^{|N(3)|^2}$ тактов, что с рациональным геометрическим множителем $\frac{d_{\mathrm{in}}(\square)\,d_{\mathrm{out}}(E)}{|N(3)| + d_{\mathrm{out}}(E)} = \frac{4\times 5}{12 + 5} = \mathbf{\frac{20}{17}}$ даёт \textbf{эталон секунды носителя без участия цезия и гравитационной постоянной~$G$}:
\begin{equation}
  \label{eq:course-carrier-second}
  \mathcal{N}_{\mathrm{sec}}(3)
  \;=\; \Bigl\lfloor \frac{20}{17}\cdot 512^{16} \Bigr\rfloor
  \;\approx\; \mathbf{2{,}623617\times 10^{43}\text{ тактов/с}} \qquad (\text{согласие с } 1/h_T \text{ лучше } 0{,}017\%!).
\end{equation}

А теперь разберёмся с главным космологическим заблуждением стандартной модели $\Lambda\mathrm{CDM}$: \textbf{почему Вселенной вовсе не $13{,}8\text{ млрд лет}$, и почему космический телескоп Джеймса Уэбба (JWST) видит на красных смещениях $z \sim 10\text{--}15$ уже полностью сформированные массивные галактики и тяжёлые ядра, которым в $13{,}8$-миллиардной модели просто не хватило бы времени вырасти?}

Посмотрим на глобальную картину вечно кипящего океана (глава~2 космологического тома) и на длинный побитово-точный прогон нашего начального слова $\mathrm{IV} = 263 = 407_8$ на 3D-решётке кубооктаэдра (запись \texttt{CE-C-04}):
\begin{itemize}[leftmargin=*]
\item \textbf{Четыре фазы глобального информационного цикла океана (по $d_{\mathrm{in}}(\square)=4$ квадрантам кольца):}
  (I)~\emph{Фаза~I (Расширение и рождение вихрей --- наша макроэпоха):} трёхгранный фазовый сдвиг $\mathrm{IV}=263$ запускает каскад рождения топологических вихрей ($b=1, n\in\mathbb{Z}$) --- барионов, ядер, звёзд, галактик и планет;
  (II)~\emph{Фаза~II (Макроскопическая релаксация на $\Tlayer$):} аннигиляция встречных вихрей перекачивает макроконтраст через вязкость $\eta/s = 41/512$ в мелкомасштабное фазовое кипение при точном сохранении всех 72 бит каждой ячейки на~$\Mlayer$;
  (III)~\emph{Фаза~III (Синхронный «лёд»):} выравнивание к синхронному вращению на вакуумном полу $\mu_{\mathrm{vac}}=1$ с нулевым относительным градиентом классов фазы;
  (IV)~\emph{Фаза~IV (Нуклеация нового пузыря):} накопление голономии Гейзенберга ($\Delta\varphi_{\mathrm{disc}} = 41$ тик) на стыке трёх доменов Киббла--Зурека снова формирует слово $\mathrm{IV}=263 = 407_8$, запуская новую Фазу~I.
\item \textbf{Четыре внутренних рубежа и две каскадные эпохи внутри самой Фазы~I:}
\end{itemize}
\begin{enumerate}[label=\textbf{\arabic*.},leftmargin=*]
\item \textbf{Рубеж~1 ($\tau = 2$, первичный всплеск) и Рубеж~2 ($\tau_1 = |{\triangle}(3)| = 8$, внутреннее фотонное горлышко Первой Эпохи):}
  после рождения 86 вихревых плакет на $\tau=2$ ($\eta_B = \frac{45}{208}\alpha_{\mathrm{fs}}^4 \approx 6{,}135\times 10^{-10}$) и первичного фотонного шара на $\tau_1 = 8$ ($1\,307$ фотонных квантов Гейзенберга) на тактах $\tau = 16\dots 64$ конденсируется первая волна из $|\square(3)| = 6$ единичных барионных вихрей.
\item \textbf{Рубеж~3 ($\tau_{\mathrm{neck}} = \phaseRingCount(3)/2 = 256$, полупериодное горлышко отцепления CMB --- ложный «Большой взрыв» модели $\Lambda\mathrm{CDM}$):}
  ровно на полупериоде кольца $\tau = 256$ пять из шести первичных вихрей сбрасывают фазу в излучение, заполняя каузальный пузырь однородной термализованной фотонной баней реликтового излучения (CMB) с вариацией плотности $< 0{,}2\%$. Но посреди этой горячей фотонной бани \textbf{выживает локализованное барионное реликтовое зерно} ($w = -1$ в узле $(-2,-3,+2)$)!
\item \textbf{Рубеж~4 ($\tau = 336\dots 512$, Вторая Эпоха после горлышка --- конденсация тяжёлых ядер $|w| = 3$ вокруг выживших зёрен):}
  сразу после прохождения горлышка $\tau_{\mathrm{neck}} = 256$ выжившее с Первой Эпохи реликтовое зерно служит готовым центром конденсации: уже на такте $\tau = 336$ буквально на соседнем ребре кубооктаэдра от него рождается тяжёлое трёхквантовое ядро ($|w| = 3 = d$) и вспыхивает кластер из $|{\triangle}(3)| = 8$ барионных узлов (а первичный нуклеосинтез даёт массовую долю гелия-4 $Y_p = \frac{1}{4}\bigl(1 - \frac{1}{4\times 13}\bigr) = \frac{51}{208} \approx \mathbf{0{,}24519}$ и дейтерий $n_{\mathrm{D}}/n_{\mathrm{H}} = \frac{5^2}{4\times 13}\alpha_{\mathrm{fs}}^2 \approx \mathbf{2{,}560\times 10^{-5}}$ в точном согласии с наблюдениями!).
\end{enumerate}

Понимаешь теперь, в чём штука? Когда стандартная космология экстраполирует разлёт видимой барионной волны Второй Эпохи назад до поверхности последнего рассеяния реликтового излучения, она упирается в \textbf{полупериодное горлышко $\tau_{\mathrm{neck}} = 256$} и принимает его за «рождение Вселенной из точки»!
\begin{itemize}[noitemsep,leftmargin=*]
\item \textbf{Видимый возраст Второй Эпохи (после горлышка $\tau_{\mathrm{neck}} = 256$)} с префактором $\frac{|N(3)|-1}{d_{\mathrm{out}}(E)|{\triangle}(3)|} = \mathbf{\frac{11}{40}}$ равен:
  \[
    N_{\mathrm{wave}}^{\mathrm{alg}}
    \;=\; \frac{11}{40}\Bigl(\frac{E_0}{\Theta_{\mathrm{CMB}}}\Bigr)^2\sqrt{\frac{\alpha_{\mathrm{fs}}}{8}}
    \;\approx\; 1{,}1398\times 10^{61}\text{ тактов}
    \qquad \bigl(t_{\mathrm{wave}} = \tau - \tau_{\mathrm{neck}} \;\longleftrightarrow\; 13{,}769\text{ млрд лет}\bigr).
  \]
\item \textbf{Настоящий полный двухэпоховый каузальный возраст нашего пузыря от старта сида $\mathrm{IV} = 263 = 407_8$} включает обе половины цикла ($\tau_{\mathrm{total}} = \tau_{\mathrm{neck}} + \tau_{\mathrm{wave}} = (c_0/c)^2\,\tau_{\mathrm{wave}}$) и вдвое больше:
  \begin{equation}
    \label{eq:course-true-two-epoch-age}
    N_{\mathrm{total}}^{\mathrm{alg}}
    \;=\; \Bigl(\frac{\caSignalSpeed}{\speedC}\Bigr)^2 N_{\mathrm{wave}}^{\mathrm{alg}}
    \;=\; 2\,N_{\mathrm{wave}}^{\mathrm{alg}}
    \;\approx\; \mathbf{2{,}2796\times 10^{61}\text{ тактов}}
    \qquad \bigl(t_{\mathrm{total}} \;\longleftrightarrow\; \mathbf{26\text{--}27{,}54\text{ млрд лет}}!\bigr).
  \end{equation}
\end{itemize}
Именно поэтому в нашей Галактике и в глубоком космосе наблюдаются объекты, чей измеренный возраст \textbf{превышает $13{,}8\text{ млрд лет}$}:
\begin{enumerate}[label=(\alph*),noitemsep]
\item \textbf{Парадокс «Звезды Мафусаила» (\texttt{HD~140283}):}
  прямой тригонометрический параллакс субгиганта \texttt{HD~140283} в гало Млечного Пути (Bond et al.) даёт возраст $\mathbf{14{,}46 \pm 0{,}8\text{ млрд лет}}$ (а ранние оценки --- до $16\text{ млрд лет}$), что старше $13{,}787$-миллиардного предела $\Lambda\mathrm{CDM}$. В двухэпоховой динамике кубооктаэдра это не ошибка измерений, а прямое следствие того, что реликтовые малометалличные звёзды гало формируются вокруг барионных зёрен Первой Эпохи ($\tau < \tau_{\mathrm{neck}} = 256$), переживших полупериодное горлышко отцепления внутри полного $27{,}54$-миллиардного цикла пузыря!
\item \textbf{Зрелые массивные галактики телескопа JWST на $z \sim 10\text{--}15$:}
  галактикам Второй Эпохи не нужно было зарождаться из бесструктурного газа с нуля за $200\text{--}300\text{ млн лет}$ --- они конденсировались непосредственно на соседних рёбрах кубооктаэдра вокруг готовых топологических зёрен Первой Эпохи (именно поэтому в 2023~году данные JWST привели к независимым оценкам полного возраста Вселенной $\sim 26{,}7\text{ млрд лет}$).
\end{enumerate}

А как вокруг молодых звёзд из протопланетного диска сами собой собираются и расслаиваются планеты? Геохимическая дифференциация планеты на ядро, мантию и кору --- это \textbf{опускание вещества в точные замкнутые оболочки кубооктаэдра}:
\begin{itemize}[noitemsep,leftmargin=*]
\item \textbf{Железное ядро ($Z = 26 = 2N_\varphi$):} сумма зарядов магния ($Z=12$) и кремния ($Z=14$), а также удвоенный алюминий ($2\times 13$) в точности равны заряду железа $Z(\mathrm{Fe}) = \mathbf{26 = 2N_\varphi(3)}$ (а удвоенный кремний $14+14 = \mathbf{28 = Z(\mathrm{Ni})}$ даёт никель!).
\item \textbf{Жидкое внешнее ядро ($\mathrm{FeS}$ / $\mathrm{NiSi}$) и глубокий карбонат мантии ($\mathrm{MgCO_3}$) на второй сфере $S_2 = 42$:}
  \[
    Z(\mathrm{FeS}) \;=\; 26 + 16 \;=\; \mathbf{42 \;=\; S_2(3)},
    \qquad
    Z(\mathrm{NiSi}) \;=\; 28 + 14 \;=\; \mathbf{42 \;=\; S_2(3)},
    \qquad
    Z(\mathrm{MgCO_3}) \;=\; 12 + 6 + 24 \;=\; \mathbf{42 \;=\; S_2(3)}!
  \]
\item \textbf{Перовскитная нижняя мантия ($\mathrm{MgSiO_3}$) и корундовые ядра суперземель ($\mathrm{Al_2O_3}$) на магической оболочке $Z = 50$:}
  $Z(\mathrm{MgSiO_3}) = 12 + 14 + 24 = \mathbf{50} = 26 + 24 = Z(\mathrm{Al_2O_3})$, а оливин верхней мантии $\mathrm{Mg_2SiO_4}$ замыкает $50 + Z(\mathrm{Ca}) = \mathbf{70 = K_P + |\square|}$.
\item \textbf{Радиогенный двигатель геодинамо ($\mathrm{U}_{92}\to\mathrm{Pb}_{82}$):} уран на третьей сфере $Z(\mathrm{U}) = S_3(3) = \mathbf{92}$ распадается в терминальный свинец $Z(\mathrm{Pb}) = 2\Delta\varphi_{\mathrm{disc}}(3) = \mathbf{82}$, выделяя ровно $92 - 82 = \mathbf{10 = Z(\mathrm{Ne})}$ протонов и $146 - 126 = \mathbf{20 = 2Z(\mathrm{Ne})}$ нейтронов тепла, поддерживающего конвекцию жидкого ядра $\mathrm{FeS}$ ($S_2=42$) и магнитосферный щит планеты!
\end{itemize}

